% !PRINTED 425
\weberpage{425}
Ist umgekehrt irgend eine Lösung $x, y$ von (8) gegeben, so ist $4n$
durch das Quadrat eines jeden gemeinschaftlichen Teilers von $x$ und $y$
teilbar. Wir bezeichnen mit $\delta$ den größten gemeinschaftlichen
Teiler von $x, y$, dessen \gesperrt{Quadrat} zugleich Teiler von $n$ ist, der also
entweder gleich dem größten gemeinschaftlichen Teiler von $x, y$ oder
gleich der Hälfte desselben ist, und setzen
\[
y = \delta y', \quad n = \delta x', % [sic]
\]
und erhalten so eine \gesperrt{eigentliche Lösung} von (7).

Wenn wir also für $-m$ alle Diskriminanten setzen, für die die Gleichung
(8) überhaupt Lösungen hat, und mit $k$ die \gesperrt{Anzahl} dieser Lösungen (mit
positivem $x$) verstehen, ferner mit $h(m)$ die Klassenzahl primitiver
Formen der Diskriminante $-m$, so ergibt sich die Formel
\weq{9}{\Sigma k\, h(m) = 2\sum \partial,}

wenn $n$ kein Quadrat ist,
\[
= 2\sum \partial + \sqrt{n} - 1,
\]
wenn $n$ ein Quadrat ist.

Hier ist noch daran zu erinnern, daß in den beiden Ausnahmefällen $m=3$,
$m=4$ unter $k$ in § 115 nur der dritte Teil oder die Hälfte der Zahl
der eigentlichen Lösungen von (8) verstanden war.

Diesem Umstand wollen wir jetzt dadurch gerecht werden, daß wir unter
$k$ die Gesamtzahl der Lösungen von (8) verstehen, aber
\[
h(3) = \tfrac{1}{3}, \quad h(4) = \tfrac{1}{2}
\]
setzen [statt $h(3)=1$, $h(4)=1$].

Es ist ferner zu bemerken, daß bei der Bildung der Summe (9) die zu dem
Wert $n'=1$ gehörigen Lösungen nicht mitgezählt sind. Dies kommt nur für
den Fall in Betracht, daß $n$ ein Quadrat ist. Es hat dann die Gleichung
(7):
\[
4 = y'^2 + m x'^2
\]
nur die folgenden eigentlichen Lösungen:
\[
\begin{array}{lllll}
m=3: & x'=1, & y'=\pm 1, & k=2, & h=\tfrac{1}{3},\\
m=4: & x'=1, & y'=0, & k=1, & h=\tfrac{1}{2},
\end{array}
\]
und es würde also, wenn man diese Ausnahme beseitigen wollte, zu (9)
noch zu addieren sein:
\[
\Sigma k h, \text{ auf } n'=1 \text{ bezogen}, = \tfrac{1}{2}+\tfrac{2}{3} = \tfrac{7}{6}.
\]
